% "La computadora de 20 vatios": the popular ("take-along") version.
% Each chapter paperback/NN.tex corresponds to the full chapter ../chapters/NN_*.tex with the same number;
% on the website the reader switches a chapter between the "Short" and "Full" versions.
\documentclass[11pt,openany]{book}
\usepackage[paperwidth=13cm,paperheight=20cm,top=1.8cm,bottom=1.8cm,inner=1.6cm,outer=1.4cm]{geometry}
\usepackage[T1]{fontenc}
\usepackage[utf8]{inputenc}
\usepackage{lmodern}
\usepackage[spanish,es-nodecimaldot,es-noshorthands,es-tabla]{babel}
\addto\extrasspanish{\spanishplainpercent}
\usepackage{amsmath,amssymb}
\usepackage{xcolor}
\usepackage[unicode,colorlinks=true,linkcolor=blue!60!black]{hyperref}
\usepackage{microtype}
\input{paperback/pencil} % minimal colour-pencil illustrations, one per chapter
\setlength{\parskip}{0.3em}
\setlength{\emergencystretch}{2em} % narrow paperback page: allow a little extra stretch
\renewcommand{\chaptermark}[1]{\markboth{\small #1}{}}
% neuron type code: first four letters of the primary mediator (as in the full edition)
\newcommand{\nt}[1]{\textsf{\textbf{#1}}}
% pointer to the full edition at the end of each chapter
\newcommand{\fullversion}[1]{\par\bigskip\noindent\small\emph{Fórmulas, modelos, ejercicios y fuentes: en la versión completa, capítulo~#1.}\normalsize}

\title{\Huge La computadora de 20 vatios\\[14pt] \Large Cómo calcula el cerebro, breve y sencillo}
\hypersetup{pdftitle={La computadora de 20 vatios. Cómo calcula el cerebro, breve y sencillo},pdfauthor={Konstantin Kladko}}
\author{\Large Konstantin~Kladko}
\date{}

\begin{document}
\frontmatter
\input{paperback/cover} % cover: a collage of the chapter drawings
% copyright page (same licence as the full edition)
\thispagestyle{empty}
\vspace*{\fill}
\noindent{\footnotesize \copyright~2026 Konstantin~Kladko.\\[4pt]
El texto y las figuras se distribuyen bajo la licencia Creative Commons Atribución-NoComercial-CompartirIgual 4.0 Internacional (CC BY-NC-SA 4.0): \url{https://creativecommons.org/licenses/by-nc-sa/4.0/}. El libro puede copiarse, usarse en la enseñanza, traducirse y adaptarse libremente con fines no comerciales, citando al autor y distribuyendo las obras derivadas bajo la misma licencia.\\[4pt]
Versión completa, fuentes y modelos: \url{https://github.com/kladkogex/20-watt-computer}.}
\clearpage
\tableofcontents
\input{paperback/00_preface}
\mainmatter
\input{paperback/01}
\input{paperback/02}
\input{paperback/03}
\input{paperback/04}
\input{paperback/05}
\input{paperback/06}
\input{paperback/07}
\input{paperback/08}
\input{paperback/09}
\input{paperback/10}
\input{paperback/11}
\input{paperback/12}
\input{paperback/13}
\input{paperback/14}
\input{paperback/15}
\input{paperback/16}
\input{paperback/17}
\input{paperback/18}
\input{paperback/19}
\input{paperback/20}
\input{paperback/21}
\input{paperback/22}
\end{document}
